% Maintained chapter-four figures. Transparent labels; all panels use explicit geometry.
% Generated by generate_ch04_geometry.py.
\newcommand{\MarsCoplanarPath}{\draw[kcurve] plot coordinates {(-3.3333333,-0.5510407) (-3.3055556,-0.5222965) (-3.2777778,-0.4937709) (-3.2500000,-0.4654674) (-3.2222222,-0.4373898) (-3.1944444,-0.4095421) (-3.1666667,-0.3819279) (-3.1388889,-0.3545514) (-3.1111111,-0.3274165) (-3.0833333,-0.3005273) (-3.0555556,-0.2738881) (-3.0277778,-0.2475030) (-3.0000000,-0.2213764) (-2.9722222,-0.1955128) (-2.9444444,-0.1699166) (-2.9166667,-0.1445923) (-2.8888889,-0.1195447) (-2.8611111,-0.0947784) (-2.8333333,-0.0702983) (-2.8055556,-0.0461093) (-2.7777778,-0.0222163) (-2.7500000,0.0013755) (-2.7222222,0.0246611) (-2.6944444,0.0476352) (-2.6666667,0.0702925) (-2.6388889,0.0926277) (-2.6111111,0.1146351) (-2.5833333,0.1363092) (-2.5555556,0.1576444) (-2.5277778,0.1786350) (-2.5000000,0.1992751) (-2.4722222,0.2195588) (-2.4444444,0.2394801) (-2.4166667,0.2590329) (-2.3888889,0.2782111) (-2.3611111,0.2970085) (-2.3333333,0.3154188) (-2.3055556,0.3334355) (-2.2777778,0.3510522) (-2.2500000,0.3682624) (-2.2222222,0.3850595) (-2.1944444,0.4014369) (-2.1666667,0.4173877) (-2.1388889,0.4329052) (-2.1111111,0.4479827) (-2.0833333,0.4626132) (-2.0555556,0.4767897) (-2.0277778,0.4905054) (-2.0000000,0.5037533) (-1.9722222,0.5165262) (-1.9444444,0.5288173) (-1.9166667,0.5406194) (-1.8888889,0.5519254) (-1.8611111,0.5627284) (-1.8333333,0.5730214) (-1.8055556,0.5827972) (-1.7777778,0.5920490) (-1.7500000,0.6007698) (-1.7222222,0.6089529) (-1.6944444,0.6165913) (-1.6666667,0.6236784) (-1.6388889,0.6302076) (-1.6111111,0.6361725) (-1.5833333,0.6415667) (-1.5555556,0.6463840) (-1.5277778,0.6506184) (-1.5000000,0.6542641) (-1.4722222,0.6573154) (-1.4444444,0.6597670) (-1.4166667,0.6616138) (-1.3888889,0.6628508) (-1.3611111,0.6634734) (-1.3333333,0.6634776) (-1.3055556,0.6628593) (-1.2777778,0.6616150) (-1.2500000,0.6597415) (-1.2222222,0.6572363) (-1.1944444,0.6540970) (-1.1666667,0.6503219) (-1.1388889,0.6459097) (-1.1111111,0.6408597) (-1.0833333,0.6351718) (-1.0555556,0.6288464) (-1.0277778,0.6218845) (-1.0000000,0.6142879) (-0.9722222,0.6060591) (-0.9444444,0.5972011) (-0.9166667,0.5877179) (-0.8888889,0.5776141) (-0.8611111,0.5668950) (-0.8333333,0.5555671) (-0.8055556,0.5436373) (-0.7777778,0.5311136) (-0.7500000,0.5180049) (-0.7222222,0.5043209) (-0.6944444,0.4900722) (-0.6666667,0.4752704) (-0.6388889,0.4599281) (-0.6111111,0.4440586) (-0.5833333,0.4276764) (-0.5555556,0.4107969) (-0.5277778,0.3934361) (-0.5000000,0.3756115) (-0.4722222,0.3573410) (-0.4444444,0.3386438) (-0.4166667,0.3195397) (-0.3888889,0.3000495) (-0.3611111,0.2801948) (-0.3333333,0.2599979) (-0.3055556,0.2394821) (-0.2777778,0.2186712) (-0.2500000,0.1975897) (-0.2222222,0.1762629) (-0.1944444,0.1547165) (-0.1666667,0.1329767) (-0.1388889,0.1110704) (-0.1111111,0.0890246) (-0.0833333,0.0668668) (-0.0555556,0.0446247) (-0.0277778,0.0223265) (0.0000000,0.0000000) (0.0277778,-0.0223265) (0.0555556,-0.0446247) (0.0833333,-0.0668668) (0.1111111,-0.0890246) (0.1388889,-0.1110704) (0.1666667,-0.1329767) (0.1944444,-0.1547165) (0.2222222,-0.1762629) (0.2500000,-0.1975897) (0.2777778,-0.2186712) (0.3055556,-0.2394821) (0.3333333,-0.2599979) (0.3611111,-0.2801948) (0.3888889,-0.3000495) (0.4166667,-0.3195397) (0.4444444,-0.3386438) (0.4722222,-0.3573410) (0.5000000,-0.3756115) (0.5277778,-0.3934361) (0.5555556,-0.4107969) (0.5833333,-0.4276764) (0.6111111,-0.4440586) (0.6388889,-0.4599281) (0.6666667,-0.4752704) (0.6944444,-0.4900722) (0.7222222,-0.5043209) (0.7500000,-0.5180049) (0.7777778,-0.5311136) (0.8055556,-0.5436373) (0.8333333,-0.5555671) (0.8611111,-0.5668950) (0.8888889,-0.5776141) (0.9166667,-0.5877179) (0.9444444,-0.5972011) (0.9722222,-0.6060591) (1.0000000,-0.6142879) (1.0277778,-0.6218845) (1.0555556,-0.6288464) (1.0833333,-0.6351718) (1.1111111,-0.6408597) (1.1388889,-0.6459097) (1.1666667,-0.6503219) (1.1944444,-0.6540970) (1.2222222,-0.6572363) (1.2500000,-0.6597415) (1.2777778,-0.6616150) (1.3055556,-0.6628593) (1.3333333,-0.6634776) (1.3611111,-0.6634734) (1.3888889,-0.6628508) (1.4166667,-0.6616138) (1.4444444,-0.6597670) (1.4722222,-0.6573154) (1.5000000,-0.6542641) (1.5277778,-0.6506184) (1.5555556,-0.6463840) (1.5833333,-0.6415667) (1.6111111,-0.6361725) (1.6388889,-0.6302076) (1.6666667,-0.6236784) (1.6944444,-0.6165913) (1.7222222,-0.6089529) (1.7500000,-0.6007698) (1.7777778,-0.5920490) (1.8055556,-0.5827972) (1.8333333,-0.5730214) (1.8611111,-0.5627284) (1.8888889,-0.5519254) (1.9166667,-0.5406194) (1.9444444,-0.5288173) (1.9722222,-0.5165262) (2.0000000,-0.5037533) (2.0277778,-0.4905054) (2.0555556,-0.4767897) (2.0833333,-0.4626132) (2.1111111,-0.4479827) (2.1388889,-0.4329052) (2.1666667,-0.4173877) (2.1944444,-0.4014369) (2.2222222,-0.3850595) (2.2500000,-0.3682624) (2.2777778,-0.3510522) (2.3055556,-0.3334355) (2.3333333,-0.3154188) (2.3611111,-0.2970085) (2.3888889,-0.2782111) (2.4166667,-0.2590329) (2.4444444,-0.2394801) (2.4722222,-0.2195588) (2.5000000,-0.1992751) (2.5277778,-0.1786350) (2.5555556,-0.1576444) (2.5833333,-0.1363092) (2.6111111,-0.1146351) (2.6388889,-0.0926277) (2.6666667,-0.0702925) (2.6944444,-0.0476352) (2.7222222,-0.0246611) (2.7500000,-0.0013755) (2.7777778,0.0222163) (2.8055556,0.0461093) (2.8333333,0.0702983) (2.8611111,0.0947784) (2.8888889,0.1195447) (2.9166667,0.1445923) (2.9444444,0.1699166) (2.9722222,0.1955128) (3.0000000,0.2213764) (3.0277778,0.2475030) (3.0555556,0.2738881) (3.0833333,0.3005273) (3.1111111,0.3274165) (3.1388889,0.3545514) (3.1666667,0.3819279) (3.1944444,0.4095421) (3.2222222,0.4373898) (3.2500000,0.4654674) (3.2777778,0.4937709) (3.3055556,0.5222965) (3.3333333,0.5510407)};}
\newcommand{\MarsInclinedPath}{\draw[TextGray,thick] plot coordinates {(-3.3295465,-2.0955284) (-3.2802682,-2.0811303) (-3.2311171,-2.0666671) (-3.1820946,-2.0521385) (-3.1332019,-2.0375442) (-3.0844402,-2.0228836) (-3.0358108,-2.0081564) (-2.9873151,-1.9933622) (-2.9389542,-1.9785006) (-2.8907295,-1.9635712) (-2.8426423,-1.9485736) (-2.7946940,-1.9335072) (-2.7468859,-1.9183719) (-2.6992194,-1.9031670) (-2.6516959,-1.8878922) (-2.6043168,-1.8725471) (-2.5570835,-1.8571312) (-2.5099975,-1.8416442) (-2.4630601,-1.8260855) (-2.4162729,-1.8104548) (-2.3696374,-1.7947516) (-2.3231551,-1.7789754) (-2.2768275,-1.7631259) (-2.2306561,-1.7472026) (-2.1846425,-1.7312051) (-2.1387883,-1.7151328) (-2.0930950,-1.6989855) (-2.0475643,-1.6827625) (-2.0021979,-1.6664635) (-1.9569973,-1.6500881) (-1.9119642,-1.6336357) (-1.8671004,-1.6171059) (-1.8224075,-1.6004983) (-1.7778873,-1.5838123) (-1.7335415,-1.5670477) (-1.6893719,-1.5502038) (-1.6453803,-1.5332802) (-1.6015685,-1.5162765) (-1.5579383,-1.4991923) (-1.5144916,-1.4820270) (-1.4712303,-1.4647802) (-1.4281562,-1.4474514) (-1.3852712,-1.4300402) (-1.3425774,-1.4125461) (-1.3000767,-1.3949687) (-1.2577710,-1.3773075) (-1.2156624,-1.3595620) (-1.1737529,-1.3417318) (-1.1320446,-1.3238164) (-1.0905395,-1.3058154) (-1.0492397,-1.2877284) (-1.0081475,-1.2695548) (-0.9672648,-1.2512942) (-0.9265939,-1.2329463) (-0.8861370,-1.2145104) (-0.8458963,-1.1959863) (-0.8058741,-1.1773734) (-0.7660726,-1.1586713) (-0.7264942,-1.1398796) (-0.6871411,-1.1209979) (-0.6480158,-1.1020257) (-0.6091205,-1.0829626) (-0.5704577,-1.0638082) (-0.5320299,-1.0445622) (-0.4938394,-1.0252240) (-0.4558888,-1.0057933) (-0.4181806,-0.9862697) (-0.3807174,-0.9666528) (-0.3435015,-0.9469423) (-0.3065358,-0.9271377) (-0.2698228,-0.9072387) (-0.2333651,-0.8872450) (-0.1971654,-0.8671562) (-0.1612265,-0.8469720) (-0.1255510,-0.8266920) (-0.0901418,-0.8063159) (-0.0550015,-0.7858435) (-0.0201331,-0.7652744) (0.0144606,-0.7446083) (0.0487768,-0.7238450) (0.0828125,-0.7029842) (0.1165648,-0.6820258) (0.1500307,-0.6609693) (0.1832073,-0.6398147) (0.2160915,-0.6185618) (0.2486803,-0.5972103) (0.2809706,-0.5757601) (0.3129592,-0.5542111) (0.3446432,-0.5325631) (0.3760191,-0.5108160) (0.4070840,-0.4889698) (0.4378346,-0.4670243) (0.4682675,-0.4449796) (0.4983796,-0.4228356) (0.5281674,-0.4005923) (0.5576278,-0.3782497) (0.5867573,-0.3558080) (0.6155525,-0.3332672) (0.6440100,-0.3106273) (0.6721263,-0.2878886) (0.6998980,-0.2650512) (0.7273216,-0.2421154) (0.7543935,-0.2190813) (0.7811103,-0.1959492) (0.8074682,-0.1727194) (0.8334638,-0.1493923) (0.8590933,-0.1259683) (0.8843532,-0.1024477) (0.9092398,-0.0788311) (0.9337493,-0.0551189) (0.9578781,-0.0313116) (0.9816224,-0.0074100) (1.0049784,0.0165855) (1.0279425,0.0406742) (1.0505107,0.0648552) (1.0726793,0.0891278) (1.0944445,0.1134912) (1.1158024,0.1379444) (1.1367491,0.1624866) (1.1572809,0.1871168) (1.1773937,0.2118338) (1.1970837,0.2366365) (1.2163470,0.2615238) (1.2351798,0.2864944) (1.2535780,0.3115470) (1.2715378,0.3366802) (1.2890553,0.3618925) (1.3061265,0.3871823) (1.3227475,0.4125482) (1.3389144,0.4379884) (1.3546232,0.4635011) (1.3698701,0.4890844) (1.3846513,0.5147366) (1.3989627,0.5404554) (1.4128005,0.5662389) (1.4261608,0.5920848) (1.4390399,0.6179909) (1.4514339,0.6439547) (1.4633391,0.6699738) (1.4747515,0.6960456) (1.4856676,0.7221674) (1.4960836,0.7483365) (1.5059959,0.7745500) (1.5154007,0.8008047) (1.5242946,0.8270978) (1.5326740,0.8534258) (1.5405354,0.8797856) (1.5478754,0.9061736) (1.5546905,0.9325863) (1.5609776,0.9590199) (1.5667332,0.9854707) (1.5719543,1.0119348) (1.5766377,1.0384079) (1.5807804,1.0648861) (1.5843795,1.0913649) (1.5874320,1.1178398) (1.5899353,1.1443064) (1.5918866,1.1707597) (1.5932834,1.1971951) (1.5941232,1.2236074) (1.5944036,1.2499916)};
\draw[DeepBlue,thick] plot coordinates {(1.5944036,1.2499916) (1.5941294,1.2760148) (1.5933054,1.3020001) (1.5919299,1.3279421) (1.5900008,1.3538355) (1.5875167,1.3796745) (1.5844759,1.4054533) (1.5808772,1.4311661) (1.5767192,1.4568067) (1.5720010,1.4823691) (1.5667215,1.5078468) (1.5608801,1.5332334) (1.5544761,1.5585223) (1.5475092,1.5837068) (1.5399791,1.6087800) (1.5318859,1.6337349) (1.5232296,1.6585643) (1.5140107,1.6832610) (1.5042297,1.7078176) (1.4938874,1.7322266) (1.4829847,1.7564803) (1.4715230,1.7805711) (1.4595036,1.8044909) (1.4469283,1.8282319) (1.4337989,1.8517860) (1.4201176,1.8751449) (1.4058869,1.8983004) (1.3911094,1.9212441) (1.3757880,1.9439675) (1.3599260,1.9664621) (1.3435267,1.9887192) (1.3265940,2.0107302) (1.3091318,2.0324862) (1.2911445,2.0539784) (1.2726366,2.0751981) (1.2536130,2.0961362) (1.2340789,2.1167838) (1.2140397,2.1371319) (1.1935012,2.1571716) (1.1724694,2.1768938) (1.1509508,2.1962896) (1.1289519,2.2153500) (1.1064797,2.2340659) (1.0835415,2.2524286) (1.0601450,2.2704290) (1.0362979,2.2880583) (1.0120084,2.3053078) (0.9872852,2.3221687) (0.9621369,2.3386324) (0.9365727,2.3546905) (0.9106020,2.3703343) (0.8842345,2.3855558) (0.8574802,2.4003466) (0.8303493,2.4146987) (0.8028525,2.4286044) (0.7750006,2.4420557) (0.7468048,2.4550453) (0.7182764,2.4675657) (0.6894271,2.4796097) (0.6602689,2.4911705) (0.6308139,2.5022413) (0.6010746,2.5128157) (0.5710636,2.5228872) (0.5407939,2.5324501) (0.5102784,2.5414984) (0.4795306,2.5500267) (0.4485639,2.5580300) (0.4173921,2.5655031) (0.3860290,2.5724416) (0.3544888,2.5788412) (0.3227855,2.5846979) (0.2909335,2.5900080) (0.2589474,2.5947682) (0.2268418,2.5989754) (0.1946313,2.6026272) (0.1623307,2.6057210) (0.1299549,2.6082551) (0.0975190,2.6102277) (0.0650378,2.6116375) (0.0325264,2.6124838) (0.0000000,2.6127660) (-0.0325264,2.6124838) (-0.0650378,2.6116375) (-0.0975190,2.6102277) (-0.1299549,2.6082551) (-0.1623307,2.6057210) (-0.1946313,2.6026272) (-0.2268418,2.5989754) (-0.2589474,2.5947682) (-0.2909335,2.5900080) (-0.3227855,2.5846979) (-0.3544888,2.5788412) (-0.3860290,2.5724416) (-0.4173921,2.5655031) (-0.4485639,2.5580300) (-0.4795306,2.5500267) (-0.5102784,2.5414984) (-0.5407939,2.5324501) (-0.5710636,2.5228872) (-0.6010746,2.5128157) (-0.6308139,2.5022413) (-0.6602689,2.4911705) (-0.6894271,2.4796097) (-0.7182764,2.4675657) (-0.7468048,2.4550453) (-0.7750006,2.4420557) (-0.8028525,2.4286044) (-0.8303493,2.4146987) (-0.8574802,2.4003466) (-0.8842345,2.3855558) (-0.9106020,2.3703343) (-0.9365727,2.3546905) (-0.9621369,2.3386324) (-0.9872852,2.3221687) (-1.0120084,2.3053078) (-1.0362979,2.2880583) (-1.0601450,2.2704290) (-1.0835415,2.2524286) (-1.1064797,2.2340659) (-1.1289519,2.2153500) (-1.1509508,2.1962896) (-1.1724694,2.1768938) (-1.1935012,2.1571716) (-1.2140397,2.1371319) (-1.2340789,2.1167838) (-1.2536130,2.0961362) (-1.2726366,2.0751981) (-1.2911445,2.0539784) (-1.3091318,2.0324862) (-1.3265940,2.0107302) (-1.3435267,1.9887192) (-1.3599260,1.9664621) (-1.3757880,1.9439675) (-1.3911094,1.9212441) (-1.4058869,1.8983004) (-1.4201176,1.8751449) (-1.4337989,1.8517860) (-1.4469283,1.8282319) (-1.4595036,1.8044909) (-1.4715230,1.7805711) (-1.4829847,1.7564803) (-1.4938874,1.7322266) (-1.5042297,1.7078176) (-1.5140107,1.6832610) (-1.5232296,1.6585643) (-1.5318859,1.6337349) (-1.5399791,1.6087800) (-1.5475092,1.5837068) (-1.5544761,1.5585223) (-1.5608801,1.5332334) (-1.5667215,1.5078468) (-1.5720010,1.4823691) (-1.5767192,1.4568067) (-1.5808772,1.4311661) (-1.5844759,1.4054533) (-1.5875167,1.3796745) (-1.5900008,1.3538355) (-1.5919299,1.3279421) (-1.5933054,1.3020001) (-1.5941294,1.2760148) (-1.5944036,1.2499916)};
\draw[TextGray,thick] plot coordinates {(-1.5944036,1.2499916) (-1.5941232,1.2236074) (-1.5932834,1.1971951) (-1.5918866,1.1707597) (-1.5899353,1.1443064) (-1.5874320,1.1178398) (-1.5843795,1.0913649) (-1.5807804,1.0648861) (-1.5766377,1.0384079) (-1.5719543,1.0119348) (-1.5667332,0.9854707) (-1.5609776,0.9590199) (-1.5546905,0.9325863) (-1.5478754,0.9061736) (-1.5405354,0.8797856) (-1.5326740,0.8534258) (-1.5242946,0.8270978) (-1.5154007,0.8008047) (-1.5059959,0.7745500) (-1.4960836,0.7483365) (-1.4856676,0.7221674) (-1.4747515,0.6960456) (-1.4633391,0.6699738) (-1.4514339,0.6439547) (-1.4390399,0.6179909) (-1.4261608,0.5920848) (-1.4128005,0.5662389) (-1.3989627,0.5404554) (-1.3846513,0.5147366) (-1.3698701,0.4890844) (-1.3546232,0.4635011) (-1.3389144,0.4379884) (-1.3227475,0.4125482) (-1.3061265,0.3871823) (-1.2890553,0.3618925) (-1.2715378,0.3366802) (-1.2535780,0.3115470) (-1.2351798,0.2864944) (-1.2163470,0.2615238) (-1.1970837,0.2366365) (-1.1773937,0.2118338) (-1.1572809,0.1871168) (-1.1367491,0.1624866) (-1.1158024,0.1379444) (-1.0944445,0.1134912) (-1.0726793,0.0891278) (-1.0505107,0.0648552) (-1.0279425,0.0406742) (-1.0049784,0.0165855) (-0.9816224,-0.0074100) (-0.9578781,-0.0313116) (-0.9337493,-0.0551189) (-0.9092398,-0.0788311) (-0.8843532,-0.1024477) (-0.8590933,-0.1259683) (-0.8334638,-0.1493923) (-0.8074682,-0.1727194) (-0.7811103,-0.1959492) (-0.7543935,-0.2190813) (-0.7273216,-0.2421154) (-0.6998980,-0.2650512) (-0.6721263,-0.2878886) (-0.6440100,-0.3106273) (-0.6155525,-0.3332672) (-0.5867573,-0.3558080) (-0.5576278,-0.3782497) (-0.5281674,-0.4005923) (-0.4983796,-0.4228356) (-0.4682675,-0.4449796) (-0.4378346,-0.4670243) (-0.4070840,-0.4889698) (-0.3760191,-0.5108160) (-0.3446432,-0.5325631) (-0.3129592,-0.5542111) (-0.2809706,-0.5757601) (-0.2486803,-0.5972103) (-0.2160915,-0.6185618) (-0.1832073,-0.6398147) (-0.1500307,-0.6609693) (-0.1165648,-0.6820258) (-0.0828125,-0.7029842) (-0.0487768,-0.7238450) (-0.0144606,-0.7446083) (0.0201331,-0.7652744) (0.0550015,-0.7858435) (0.0901418,-0.8063159) (0.1255510,-0.8266920) (0.1612265,-0.8469720) (0.1971654,-0.8671562) (0.2333651,-0.8872450) (0.2698228,-0.9072387) (0.3065358,-0.9271377) (0.3435015,-0.9469423) (0.3807174,-0.9666528) (0.4181806,-0.9862697) (0.4558888,-1.0057933) (0.4938394,-1.0252240) (0.5320299,-1.0445622) (0.5704577,-1.0638082) (0.6091205,-1.0829626) (0.6480158,-1.1020257) (0.6871411,-1.1209979) (0.7264942,-1.1398796) (0.7660726,-1.1586713) (0.8058741,-1.1773734) (0.8458963,-1.1959863) (0.8861370,-1.2145104) (0.9265939,-1.2329463) (0.9672648,-1.2512942) (1.0081475,-1.2695548) (1.0492397,-1.2877284) (1.0905395,-1.3058154) (1.1320446,-1.3238164) (1.1737529,-1.3417318) (1.2156624,-1.3595620) (1.2577710,-1.3773075) (1.3000767,-1.3949687) (1.3425774,-1.4125461) (1.3852712,-1.4300402) (1.4281562,-1.4474514) (1.4712303,-1.4647802) (1.5144916,-1.4820270) (1.5579383,-1.4991923) (1.6015685,-1.5162765) (1.6453803,-1.5332802) (1.6893719,-1.5502038) (1.7335415,-1.5670477) (1.7778873,-1.5838123) (1.8224075,-1.6004983) (1.8671004,-1.6171059) (1.9119642,-1.6336357) (1.9569973,-1.6500881) (2.0021979,-1.6664635) (2.0475643,-1.6827625) (2.0930950,-1.6989855) (2.1387883,-1.7151328) (2.1846425,-1.7312051) (2.2306561,-1.7472026) (2.2768275,-1.7631259) (2.3231551,-1.7789754) (2.3696374,-1.7947516) (2.4162729,-1.8104548) (2.4630601,-1.8260855) (2.5099975,-1.8416442) (2.5570835,-1.8571312) (2.6043168,-1.8725471) (2.6516959,-1.8878922) (2.6992194,-1.9031670) (2.7468859,-1.9183719) (2.7946940,-1.9335072) (2.8426423,-1.9485736) (2.8907295,-1.9635712) (2.9389542,-1.9785006) (2.9873151,-1.9933622) (3.0358108,-2.0081564) (3.0844402,-2.0228836) (3.1332019,-2.0375442) (3.1820946,-2.0521385) (3.2311171,-2.0666671) (3.2802682,-2.0811303) (3.3295465,-2.0955284)};
\draw[->,DeepBlue,thick] plot coordinates {(-0.8732643,-1.2086051) (-0.8296331,-1.1884459) (-0.7862609,-1.1681816) (-0.7431505,-1.1478119) (-0.7003050,-1.1273361) (-0.6577273,-1.1067537) (-0.6154204,-1.0860642) (-0.5733875,-1.0652670) (-0.5316316,-1.0443616)};
\draw[->,DeepBlue,thick] plot coordinates {(0.6894671,2.4795935) (0.6573750,2.4922850) (0.6249252,2.5043828) (0.5921344,2.5158783) (0.5590194,2.5267631) (0.5255976,2.5370295) (0.4918864,2.5466700) (0.4579039,2.5556774) (0.4236680,2.5640451)};
\draw[->,DeepBlue,thick] plot coordinates {(0.8732643,-1.2086051) (0.9171514,-1.2286600) (0.9612918,-1.2486111) (1.0056825,-1.2684589) (1.0503207,-1.2882040) (1.0952037,-1.3078471) (1.1403288,-1.3273886) (1.1856933,-1.3468291) (1.2312945,-1.3661693)};}

\newcommand{\ReferenceFramesDiagram}{%
\begin{center}\begin{minipage}{\linewidth}\centering\begin{tikzpicture}[>=Latex,scale=.88]
\draw[kaxis] (-.3,0)--(7,0) node[right]{$S:x$};
\draw[kaxis] (2,1.6)--(7,1.6) node[right]{$S':x'$};
\coordinate (P) at (5.7,2.4);
\draw[kguide] (2,-.6)--(2,1.6) (P)--(5.7,-.6);
\fill (0,0) circle(1.5pt) (2,1.6) circle(1.5pt) (P) circle(1.5pt);
\node[below left] at (0,0){$O$};\node[above left] at (2,1.6){$O'$};
\node[above] at (P){$P$};
\draw[<->] (0,-.65)--(2,-.65) node[midway,below]{$x_{O'}$};
\draw[<->,DeepBlue] (0,-1.25)--(5.7,-1.25) node[midway,below]{$x$};
\draw[<->,DeepBlue] (2,.8)--(5.7,.8) node[midway,above]{$x'$};
\draw[vector] (2.6,2.35)--(3.9,2.35) node[midway,above]{$u$};
\draw (5.7,0)--(5.7,.12) (5.7,1.6)--(5.7,1.72);
\end{tikzpicture}\figcap{图2.19：同一点的两个横向投影坐标满足 $x=x_{O'}+x'$；轴的上下错开不表示额外横向距离。}\end{minipage}\end{center}}

\newcommand{\AircraftWindDiagram}{%
\begin{center}\begin{minipage}{\linewidth}\centering\begin{tikzpicture}[>=Latex,scale=.88]
\begin{scope}
\draw[kaxis] (0,0)--(4,0) node[right]{东};\draw[kaxis] (0,0)--(0,2.4) node[above]{北};
\draw[vector] (0,0)--(3,0) node[midway,below]{$\vect V$};
\draw[vector] (3,0)--(3,1.5) node[midway,right]{$\vect u$};
\draw[vector] (0,0)--(3,1.5) node[midway,above left]{$\vect v_{\text{地}}$};
\node at (1.7,-.85){(a) 机头正东};
\end{scope}
\begin{scope}[xshift=7cm]
\draw[kaxis] (0,0)--(3.8,0) node[right]{东};\draw[kaxis] (0,0)--(0,2.4) node[above]{北};
\draw[vector] (0,0)--(2.598,-1.5) node[midway,below left]{$\vect V$};
\draw[vector] (2.598,-1.5)--(2.598,0) node[midway,right]{$\vect u$};
\draw[vector] (0,0)--(2.598,0) node[midway,above]{$\vect v_{\text{地}}$};
\draw (.8,0) arc(0:-30:.8);\node at (1.15,-.3){$\theta$};
\node at (1.5,-2.1){(b) 航迹正东};
\end{scope}
\end{tikzpicture}\figcap{例2.81解图：风速向北；两图均按 $\vect v_{\text{地}}=\vect V+\vect u$ 首尾相接。}\end{minipage}\end{center}}

\newcommand{\RiverCrossingDiagram}{%
\begin{center}\begin{minipage}{\linewidth}\centering\begin{tikzpicture}[>=Latex,scale=.9]
\begin{scope}
\fill[Blue!5] (-.5,0) rectangle (3.9,3);
\draw[thick] (-.5,0)--(3.9,0) (-.5,3)--(3.9,3);
\coordinate (A) at (0,0);\draw[kguide] (A)--(2.25,3);
\fill (A) circle(1.5pt);\node[below] at (A){$A$};
\draw[vector] (A)--(0,1.6) node[midway,left]{$\vect v_1$};
\draw[vector] (0,1.6)--(1.2,1.6) node[midway,above]{$\vect v_2$};
\draw[vector] (A)--(1.2,1.6);
\node[right] at (1.55,.65){$\vect v_{\text{岸}}$};\draw[thin] (1.55,.78)--(.82,1.02);
\draw[<->] (3.5,0)--(3.5,3) node[midway,left]{$d$};
\node at (1.7,-.8){(a) 最短时间的实际航迹};
\end{scope}
\begin{scope}[xshift=7cm,yshift=.35cm]
\coordinate (W) at (2.5,0);\coordinate (T) at (1.6,1.2);
\draw[kaxis] (-.3,0)--(4.4,0) node[right]{$v_x$};
\draw[kaxis] (0,-1.75)--(0,2.3) node[above]{$v_y$};
\draw (W) circle(1.5);\draw[kguide] (W)--(T);
\draw[vector] (0,0)--(T);\fill (W) circle(1.5pt) (T) circle(1.5pt);
\node[below left] at (0,0){$O$};\node[below] at (W){$W$};\node[above] at (T){$T$};
\node[above right] at (2.2,.65){$v_1$};
\draw[<->] (0,-1.9)--(2.5,-1.9) node[midway,below]{$v_2$};
\draw (.4,0) arc(0:36.87:.4);\node at (.67,.2){$\psi$};
\node at (1.8,-2.65){(b) $v_1<v_2$ 的速度圆};
\end{scope}
\end{tikzpicture}\figcap{例2.82解图：左图出发点在岸上，航迹平行于合速度；右图 $WT\perp OT$，切线给出最小位移方向。}\end{minipage}\end{center}}

\newcommand{\AircraftEqualSpeedDiagram}{%
\begin{center}\begin{minipage}{\linewidth}\centering\begin{tikzpicture}[>=Latex,scale=1]
\foreach \s/\shift/\txt in {1/0/东,-1/6/西}{
\begin{scope}[xshift=\shift cm]
\coordinate (O) at (0,0);\coordinate (G) at (0,3.6);
\coordinate (A) at ({\s*1.74284},3.15);\coordinate (F) at ($(A)!.5!(G)$);
\draw[vector] (O)--(A);\draw[vector] (A)--(G);\draw[vector] (O)--(G);
\draw[kguide] (O)--(F);
\node[below] at (O){$O$};\node[above] at (G){$G$};
\ifnum\s>0 \node[right] at (A){$A$};\node[left] at (0,1.45){$60$};
\else \node[left] at (A){$A$};\node[right] at (0,1.45){$60$};\fi
\node[above] at (F){$F$};\node at ({\s*1.35},1.6){$60$};
\draw (90:1.4) arc[start angle=90,end angle={90-\s*28.955},radius=1.4];
\node at ({\s*.61},1.77){$\theta$};
\node at ({\s*.9},3.9){$30$};
\node at (0,-.65){机头北偏\txt};
\end{scope}}
\node at (3,2.1){正北 $\uparrow$};
\end{tikzpicture}\figcap{例2.83解图：$OA$ 为空速，$AG$ 为风速，$OG$ 为地速；$F$ 为风速边中点，速度单位均为 $\mathrm{m/s}$。}\end{minipage}\end{center}}

\newcommand{\SmokeWindDiagram}{%
\begin{center}\begin{minipage}{\linewidth}\centering\begin{tikzpicture}[>=Latex,scale=.95]
\foreach \vv/\shift/\txt in {1/0/第一次,3/6/第二次}{
\begin{scope}[xshift=\shift cm]
\draw[kaxis] (-.3,0)--(3.3,0) node[right]{东};\draw[kaxis] (0,-1.6)--(0,2.6) node[above]{北};
\draw[vector] (0,0)--(2,-1) node[midway,below left]{$\vect w$};
\draw[vector] (2,-1)--(2,{\vv-1});
\node[right] at (2.1,{.5*\vv-1}){$-\vect v_{\text{船}}$};
\draw[vector] (0,0)--(2,{\vv-1});
\node[above left] at (1,{.5*(\vv-1)+.2}){$\vect v_{\text{烟/船}}$};
\node at (1.4,-2.1){\txt};
\end{scope}}
\end{tikzpicture}\figcap{例2.84解图：同一风速加上两种反向船速，分别得到向东、向东北的烟对船速度。}\end{minipage}\end{center}}

\newcommand{\InterceptTriangleDiagram}{%
\begin{center}\begin{minipage}{\linewidth}\centering\begin{tikzpicture}[>=Latex,x=.045cm,y=.045cm]
\coordinate (A) at (0,0);\coordinate (B) at (-93.9693,34.2020);
\coordinate (C1) at (-17.3065,34.2020);\coordinate (C2) at (79.9527,34.2020);
\draw[thick] (A)--(B);\draw[vector] (B)--(C2);
\draw[vector] (A)--(C1);\draw[vector] (A)--(C2);
\foreach \p/\lab/\pos in {A/A/below,B/B/above left,C1/C_1/above,C2/C_2/above right}{\fill (\p) circle(1.6pt);\node[\pos] at (\p){$\lab$};}
\node[left] at (-50,12){$L$};\node[above] at (30,38){敌船沿 $+x$};
\node[left] at (-8,20){$v_2t_1$};\node[below right] at (42,17){$v_2t_2$};
\draw (B)++(0:22) arc(0:-20:22);\node at (-64,29){$\alpha$};
\draw (23.16018:14) arc(23.16018:160:14);\node[below] at (0,-12){$\beta_1=\angle BAC_1,\quad\beta_2=\angle BAC_2$};
\end{tikzpicture}\figcap{例2.85解图：按文中数值例得到的两个未来交点；发射方向由各自的 $AC_i$ 决定，而不只由正弦值决定。}\end{minipage}\end{center}}

\newcommand{\BombInterceptDiagram}{%
\begin{center}\begin{minipage}{\linewidth}\centering\begin{tikzpicture}[>=Latex,scale=.92]
\draw[kaxis] (-5,0)--(3.8,0) node[right]{$x$};
\coordinate (B) at (-4,3);\coordinate (Q) at (-2,2.25);
\draw[kcurve,domain=0:1,samples=60,variable=\t] plot({-4+4*\t},{3*(1-\t*\t)});
\draw[kguide] (B)--(-4,0);\draw[vector] (B)--(-2.8,3) node[midway,above]{$v_1$};
\draw[thick,dashed] (-2,0)--(Q) (2,0)--(Q);
\foreach \x/\lab/\pos in {0/0/below,-2/-D/below,2/+D/below}{\fill (\x,0) circle(1.6pt);\node[\pos] at (\x,0){$\lab$};}
\node[above left] at (B){$(-L,H)$};\fill (Q) circle(1.6pt);\node[above right] at (Q){$Q$};
\draw[<->] (-4.65,0)--(-4.65,3) node[midway,left]{$H$};
\draw[<->] (-4,-.8)--(0,-.8) node[midway,below]{$L$};
\node at (1.6,3){两炮位分别讨论};
\end{tikzpicture}\figcap{例2.86解图：蓝线是炸弹轨迹；两条虚线仅连接炮位与相遇点，不是炮弹的抛物线轨迹。}\end{minipage}\end{center}}

\newcommand{\ThreeBoatsDiagram}{%
\begin{center}\begin{minipage}{\linewidth}\centering\begin{tikzpicture}[>=Latex,scale=.95]
\begin{scope}
\draw[kguide] (0,0)--(4.8,0);
\foreach \x/\lab in {0/A,1.8/B,3.6/C}{\fill (\x,0) circle(1.5pt);\node[below] at (\x,0){$\lab$};}
\draw[vector] (0,0)--(.6,0);\node[above] at (.3,0){$u$};
\draw[vector] (1.8,0)--(1.8,1.2) node[above]{$2u$};
\draw[vector] (3.6,0)--(4.8,1.2) node[above]{$2\sqrt2u$};
\draw[<->] (0,-.65)--(1.8,-.65) node[midway,below]{$a$};
\draw[<->] (1.8,-.65)--(3.6,-.65) node[midway,below]{$a$};
\node at (2,2.1){(a) $t=0$，速度同比例};
\end{scope}
\begin{scope}[xshift=7cm]
\coordinate (M) at (2.7,2.2);\coordinate (A) at (0,0);
\draw[kguide] (A)--(2.7,0)--(M);\draw[thick,->] (M)--(A);
\fill (M) circle(1.5pt) (A) circle(1.5pt);
\node[above] at (M){$M(t)$};\node[left] at (A){$A(t)$};
\node[below] at (1.35,0){$3a/2$};\node[right] at (2.7,1){$2ut$};
\draw[vector] (A)--(0,-1.2) node[below]{$\vect v_{A/M}=(0,-2u)$};
\node[above left] at (1.5,1.3){$\vect r_{A/M}$};
\node at (1.5,3.05){(b) 随 $M$ 平移的坐标};
\end{scope}
\end{tikzpicture}\figcap{例2.87解图：$M$ 与 $A$ 的横向速度相同，$MA$ 的水平分量固定、竖直分量随时间增加。}\end{minipage}\end{center}}

\newcommand{\PursuitDiagram}{%
\begin{center}\begin{minipage}{\linewidth}\centering\begin{tikzpicture}[>=Latex,scale=1.25]
\draw[thick] (-.3,0)--(3.5,0);
\draw[kguide] (0,3)--(3.1,3);
\draw[kcurve,domain=1:0,samples=100,variable=\u] plot({3*(2/3-\u+\u^3/3)},{3*(1-\u^2)});
\coordinate (D) at (.625,2.25);\coordinate (P) at (1.1875,3);
\draw[thick] (D)--(P);\draw[kguide] (D)--++(1.3,0);
\draw[vector] (D)--++(.27,.36);\draw[vector] (P)--++(.5,0) node[midway,above]{$v_0$};
\fill (D) circle(1.5pt) (P) circle(1.5pt);
\node[left] at (D){$D$};\node[above left] at (P){$P$};
\node[left] at (.86,2.7){$r$};
\draw (D)++(.38,0) arc(0:53.1301:.38);\node at (1.18,2.52){$\theta$};
\draw[vector] (0,0)--(0,.7) node[left]{$v_1$};
\draw[vector] (0,0)--(.65,0) node[below]{$a_0$};
\draw[<->] (-.65,0)--(-.65,3) node[midway,left]{$H$};
\node[below left] at (0,0){$A$};\node[above right] at (2,3){命中};
\end{tikzpicture}\figcap{例2.88解图：取 $v_1=2v_0$ 展示追击曲线；中途导弹 $D$ 的切线沿同一时刻飞机 $P$ 的视线，初始竖直、末端水平。}\end{minipage}\end{center}}

\newcommand{\SlidingRodDiagram}{%
\begin{center}\begin{minipage}{\linewidth}\centering\begin{tikzpicture}[>=Latex,scale=.88]
\begin{scope}
\coordinate (O) at (0,0);\coordinate (A) at (0,3.6);\coordinate (B) at (4,0);\coordinate (P) at (1.52,2.232);
\draw[thick] (0,4.1)--(O)--(4.9,0);\draw[kcurve] (A)--(B);
\foreach \p/\pos in {A/left,B/below,P/above right}{\fill (\p) circle(1.5pt);\node[\pos] at (\p){$\p$};}
\draw[vector] (A)--++(0,-.8) node[midway,left]{$v_A$};\draw[vector] (B)--++(.8,0) node[above]{$v_B$};
\draw (0,2.9) arc(-90:-41.987:.7);\node at (.42,2.7){$\theta$};
\node[above right] at (.75,3.25){$aL$};\node[above right] at (2.7,1.4){$(1-a)L$};
\end{scope}
\begin{scope}[xshift=7cm]
\draw[kaxis] (0,0)--(3.3,0) node[right]{$x$};\draw[kaxis] (0,0)--(0,4.1) node[above]{$y$};
\draw[kcurve,domain=0:90,samples=60] plot({.38*sqrt(28.96)*sin(\x)},{.62*sqrt(28.96)*cos(\x)});
\coordinate (P) at (1.52,2.232);\fill (P) circle(1.5pt);
\draw[vector] (P)--++(.342,-.62);\node[above right] at (P){$P$};
\node[right] at (2,1.55){$\vect v_P$};
\node[above left] at (0,3.337){$(1-a)L$};\node[below] at (2.045,0){$aL$};
\end{scope}
\end{tikzpicture}\figcap{例2.89解图：左图杆长固定、端点接触墙与地面；右图使用同一 $P$，速度沿四分之一椭圆的切线。}\end{minipage}\end{center}}

\newcommand{\BoatTowDiagram}{%
\begin{center}\begin{minipage}{\linewidth}\centering\begin{tikzpicture}[>=Latex,scale=.93]
\coordinate (O) at (0,2.5);\coordinate (B) at (4,0);
\draw[thick] (-3.8,2.3)--(-.15,2.3)--(-.15,-.4);
\fill[Blue!5] (0,-.7) rectangle (5.3,0);\draw[thin,Blue] (0,0)--(5.3,0);
\draw[thick] (-3.2,2.5) rectangle (-2.2,3);\draw (-3,2.4) circle(.1) (-2.4,2.4) circle(.1);
\draw[kcurve] (-2.2,2.5)--(O)--(B);
\fill (O) circle(1.8pt);\node[above right] at (O){$O$};
\draw[thick] (B)--(5.1,0)--(4.9,-.4)--(4.2,-.4)--cycle;
\fill (B) circle(1.5pt);\node[above right] at (B){$B$};
\draw[vector] (-2.3,3.45)--(-3.4,3.45) node[midway,above]{$v$};
\draw[vector] (4.9,-1)--(3.9,-1) node[midway,below]{$u$};
\node[above] at (-1.2,2.5){$l_c$};\node[above right] at (1.9,1.32){$l$};
\draw[<->] (-4.25,0)--(-4.25,2.5) node[midway,left]{$h$};
\draw[kguide] (-4.25,0)--(0,0) (-4.25,2.5)--(-3.5,2.5);
\draw (B)++(148:.9) arc(148:180:.9);\node at (2.82,.31){$\theta$};
\draw[<->] (0,-.85)--(4,-.85) node[midway,below]{$x$};
\end{tikzpicture}\figcap{例2.90解图：$O$ 表示半径忽略的定滑轮转向点；汽车侧增长 $\dot l_c=v$，船侧缩短 $\dot l=-v$。}\end{minipage}\end{center}}

\newcommand{\MovingConstraintDiagram}{%
\begin{center}\begin{minipage}{\linewidth}\centering\begin{tikzpicture}[>=Latex,scale=.9]
\begin{scope}
\draw[thick] (0,0)--(4,0)--(0,2)--cycle;
\coordinate (P) at (2.2,.9);\draw[very thick] (P)--(2.2,3.1);
\foreach \x in {2.06,2.34}{\draw[thick] (\x,2)--(\x,2.6);}
\draw (1.6,2)--(2.06,2) (2.34,2)--(2.8,2);
\fill (P) circle(1.5pt);\node[above right=3pt] at (P){$P$};
\draw[vector] (.6,-.5)--(1.7,-.5) node[midway,below]{$v_0$};
\draw[vector] (2.7,2.65)--(2.7,3.5) node[right]{$v$};
\draw (3.2,0) arc(180:153.435:.8);\node at (2.95,.24){$\theta$};
\node at (1.8,3.9){(a) 斜面与固定导向};
\end{scope}
\begin{scope}[xshift=7cm]
\coordinate (C) at (0,0);\coordinate (P) at (1.2,1.6);
\draw[thick] (-2,0)--(2,0) arc(0:180:2);
\draw[kguide] (C)--(0,2.2) (C)--(P) (P)--(1.2,0);
\draw[very thick] (P)--(1.2,3.2);
\foreach \x in {1.06,1.34}{\draw[thick] (\x,2.4)--(\x,3);}
\draw (.6,2.4)--(1.06,2.4) (1.34,2.4)--(1.8,2.4);
\fill (P) circle(1.5pt) (C) circle(1.5pt);
\node[below left] at (C){$C=(X,0)$};\node[above right=3pt] at (P){$P$};
\node at (.7,.46){$R$};\node[below] at (1.2,0){$x_0$};
\draw (0,.75) arc(90:53.1301:.75);\node at (.29,1.13){$\theta$};
\draw[vector] (1.8,2.7)--(1.8,3.5) node[right]{$v_y$};
\draw[vector] (.5,-.8)--(1.6,-.8) node[midway,below]{$v_0$};
\node at (0,3.9){(b) 半圆柱与固定导向};
\end{scope}
\end{tikzpicture}\figcap{例2.91—2.92解图：两侧固定导向只允许杆竖直平移；接触点由同一斜面方程或圆方程确定。}\end{minipage}\end{center}}

\newcommand{\RodPulleyDiagram}{%
\begin{center}\begin{minipage}{\linewidth}\centering\begin{tikzpicture}[>=Latex,scale=1]
\coordinate (O) at (0,0);\coordinate (B) at (0,3.4);\coordinate (A) at (2.3,1.2);\coordinate (C) at (-2.2,3.4);
\draw[very thick] (O)--(A) node[midway,below=5pt]{$R$};
\draw[kcurve] (A)--(B)--(C)--(-2.2,.65);
\draw[thick] (-2.55,0) rectangle (-1.85,.65);\node at (-2.2,.32){$M$};
\draw[kguide] (O)--(B);
\foreach \p/\pos in {O/below,B/above,C/above,A/right}{\fill (\p) circle(1.7pt);\node[\pos] at (\p){$\p$};}
\draw[<->] (-.65,0)--(-.65,3.4) node[midway,left]{$H$};
\draw (B)++(-90:.8) arc(-90:-43.727:.8);\node at (.42,2.4){$\alpha$};
\draw[vector] (A)--++(-.48,.92) node[above]{$\vect v_A$};
\draw[->,DeepBlue] (35:.75) arc(35:120:.75);\node[above] at (.3,1.2){$\omega$};
\draw[vector] (-2.85,.6)--(-2.85,-.3) node[below]{$v_M$};
\node[left] at (.83,2.0){$l$};
\end{tikzpicture}\figcap{例2.93解图：零半径转向点 $B,C$ 与绳折点重合；正角速度下 $AB$ 缩短而物块下降。}\end{minipage}\end{center}}

\newcommand{\EulerStepsDiagram}{%
\begin{center}\begin{minipage}{\linewidth}\centering\begin{tikzpicture}[>=Latex,scale=.87]
\begin{scope}
\draw[kaxis] (-.2,0)--(2,0) node[right]{$x$};\draw[kaxis] (0,-.2)--(0,2) node[above]{$y$};
\draw[vector] (0,0)--(35:2) node[above right]{$N$};
\draw[->] (.85,0) arc(0:35:.85);\node at (17.5:1.16){$\alpha$};
\node[below left] at (0,0){$O$};\node at (.7,-.85){(1) 绕 $z$};\node at (.7,-1.35){从 $+z$ 端看};
\end{scope}
\begin{scope}[xshift=5cm]
\draw[TextGray,thick] (-1.6,0)--(1.7,0) node[right]{$xy$};
\draw[DeepBlue,thick] (215:1.7)--(35:1.7) node[right]{$XY$};
\draw[kaxis] (0,0)--(90:2.1) node[above]{$z$};
\draw[vector] (0,0)--(125:2.1) node[above left]{$Z$};
\draw[->] (90:1.05) arc(90:125:1.05);\node at (107.5:1.45){$\beta$};
\draw (0,0) circle(.085);\fill (0,0) circle(.02);
\node[right] at (.12,-.4){$N$ 向外};
\node at (0,-1.4){(2) 绕 $N$};\node at (0,-1.9){从 $+N$ 端看};
\end{scope}
\begin{scope}[xshift=9.8cm]
\draw[kguide] (-.2,0)--(2,0) node[right]{$N$};
\draw[vector] (0,0)--(35:2) node[above right]{$X$};
\draw[vector] (0,0)--(125:1.9) node[above left]{$Y$};
\draw[->] (.85,0) arc(0:35:.85);\node at (17.5:1.16){$\gamma$};
\node[below left] at (0,0){$O$};\node at (.5,-.85){(3) 绕 $Z$};\node at (.5,-1.35){从 $+Z$ 端看};
\end{scope}
\end{tikzpicture}\figcap{图2.20：右手正转的三步欧拉角；第二图中 $XY$ 截线转过 $\beta$，其法线 $Z$ 同转过 $\beta$，二者始终垂直。}\end{minipage}\end{center}}

\newcommand{\RigidMotionTypesDiagram}{%
\begin{center}\begin{minipage}{\linewidth}\centering\begin{tikzpicture}[>=Latex,scale=.9]
\begin{scope}
\foreach \dy in {0,1}{\draw[kguide,domain=0:1,samples=45,variable=\t] plot({3*\t},{\dy+.9*sin(180*\t)});}
\draw[thick] (-.3,-.25) rectangle (.3,1.25);
\draw[thick] (2.7,-.25) rectangle (3.3,1.25);
\foreach \x in {0,3}{\foreach \y in {0,1}{\fill (\x,\y) circle(1.5pt);}}
\draw[vector] (0,0)--(.6,.565);\draw[vector] (0,1)--(.6,1.565);
\node[left] at (-.45,0){$P_1$};\node[left] at (-.45,1){$Q_1$};
\node[right] at (3.45,0){$P_2$};\node[right] at (3.45,1){$Q_2$};
\node at (1.5,2.4){(a) 曲线平动};
\node at (1.5,-1){两姿态的对应边平行};
\end{scope}
\begin{scope}[xshift=8cm]
\draw[kguide] (0,-1.2)--(0,1.2);
\draw[thick] (0,-2)--(0,-1.2) node[pos=0,below]{$B$};
\draw[kaxis] (0,1.2)--(0,2.4) node[above]{$A$};
\draw[thick,dashed] (2,0) arc(0:180:2 and 1.2);
\draw[thick] (-2,0) arc(180:360:2 and 1.2);
\coordinate (P) at ({2*cos(315)},{1.2*sin(315)});
\draw[thick] (0,0)--(P);\draw[kguide] (0,0)--(2,0);
\draw[vector] (P)--++(.78,.468);
\draw (.7,0) arc(0:-45:.7 and .42);
\node at (1,-.19){$\theta$};\node at (.46,-.62){$R$};
\fill (P) circle(1.5pt);\node[below] at (P){$P$};\node[left] at (0,0){$O$};
\node[above right] at (2.2,-.38){$\vect v_P$};
\node at (0,-2.7){(b) 定轴转动的圆平面投影};
\end{scope}
\end{tikzpicture}\figcap{图2.21—图2.22：平动可以沿曲线；定轴转动中 $OP=R$，点 $P$ 的切向速度为 $|\omega|R$。}\end{minipage}\end{center}}

\newcommand{\RigidBarCasesDiagram}{%
\begin{center}\begin{minipage}{\linewidth}\centering\begin{tikzpicture}[>=Latex,scale=.95]
\begin{scope}[xscale=1.5]
\draw[thick] (0,0)--(2.4,0);\draw[kguide] (2.4,0)--(4.8,0);
\draw[vector] (0,0)--(0,1.2) node[above]{$v_0$};\draw[vector] (2.4,0)--(2.4,.6) node[above]{$v_0/2$};
\foreach \x/\lab in {0/A,2.4/B,4.8/I}{\fill (\x,0) circle(1.5pt);\node[below] at (\x,0){$\lab$};}
\fill (1.2,0) circle(1.5pt);\node[above] at (1.2,0){$G$};
\draw[<->] (0,-.7)--(2.4,-.7) node[midway,below]{$l$};\draw[<->] (2.4,-.7)--(4.8,-.7) node[midway,below]{$l$};
\node at (2.4,2){(1) $x_G=l/2,\ x_I=2l$};
\end{scope}
\begin{scope}[yshift=-4.1cm,xscale=1.5]
\draw[thick] (0,0)--(2.4,0);\draw[vector] (0,0)--(0,1.2) node[above]{$v_0$};\draw[vector] (2.4,0)--(2.4,-1.2) node[below]{$-v_0$};
\node[left] at (0,0){$A$};\node[right] at (2.4,0){$B$};\fill (1.2,0) circle(1.6pt);\node[above] at (1.2,0){$G=I$};
\node at (1.2,2){(2) 等质量};
\end{scope}
\begin{scope}[xshift=7cm,yshift=-4.1cm]
\draw[thick] (0,0)--(3.6,0);\draw[vector] (0,0)--(0,1.2) node[above]{$v_0$};\draw[vector] (3.6,0)--(3.6,-1.2) node[below]{$-v_0$};
\node[left] at (0,0){$A$};\node[right] at (3.6,0){$B$};
\fill (1.2,0) circle(1.6pt) (1.8,0) circle(1.6pt);\node[above] at (1.2,0){$G$};\node[below] at (1.8,0){$I$};
\node at (1.8,2){(3) $m_A=2m_B$};
\draw[<->] (0,-.65)--(1.2,-.65) node[midway,below]{$l/3$};
\end{scope}
\end{tikzpicture}\figcap{例2.94解图：三问分别标出实际速度、质心 $G$ 与瞬心 $I$；质量改变不改变给定速度场的瞬心。}\end{minipage}\end{center}}

\newcommand{\RollingBallDiagram}{%
\begin{center}\begin{minipage}{\linewidth}\centering\begin{tikzpicture}[>=Latex,scale=1]
\draw[thick] (-2,0)--(4,0);\draw[thick] (0,1) circle(1);
\fill (0,1) circle(1.5pt) (0,0) circle(1.5pt);\node[below left] at (0,1){$C$};\node[below left] at (0,0){$P$};
\draw (0,1)--(-.707,1.707);\node[left] at (-.58,1.28){$R$};
\draw[vector] (0,1)--(2.6,1);\node[above] at (2,1){$v_0$};
\draw[->,DeepBlue,thick] (40:1.5)++(0,1) arc(40:145:1.5);
\node[above] at (0,2.65){$\omega_0>0$（逆时针）};
\draw[vector] (0,-.65)--(2.5,-.65) node[midway,below]{$v_P=v_0+R\omega_0$};
\end{tikzpicture}\figcap{例2.95解图：接触点的平动速度和转动速度初始都向右；纯滚动要求它们相消。}\end{minipage}\end{center}}

% Knife local coordinates in metres: handle end -0.15, COM 0, tip +0.10.
\newcommand{\KnifeOutline}{%
\draw[thick,fill=TextGray!16] (-.15,-.012) rectangle (-.045,.012);
\draw[thick,fill=Blue!8] (-.045,-.012)--(.10,0)--(-.045,.012)--cycle;}

\newcommand{\KnifeFlightDiagram}{%
\begin{center}\begin{minipage}{\linewidth}\centering\begin{tikzpicture}[>=Latex]
\begin{scope}[x=2.1cm,y=2.1cm]
\draw[thick] (5,-.1)--(5,1.1);\node[right] at (5,.7){木靶};
\draw[kguide,domain=0:1,samples=100,variable=\t] plot({4.9*\t},{4.9*tan(30)*\t*(1-\t)});
\foreach \t/\n in {0/1,.25/2,.5/3,.75/4,1/5}{
\begin{scope}[shift={({4.9*\t},{4.9*tan(30)*\t*(1-\t)})},rotate={60+300*\t}]
\KnifeOutline\fill (0,0) circle(.7pt);
\end{scope}
\node[above] at ({4.9*\t},{4.9*tan(30)*\t*(1-\t)+.17}){\n};}
\draw[<->] (0,-.35)--(5,-.35) node[midway,below]{$s=5.0\,\mathrm m$};
\node at (2.5,-.8){(a) 质心抛物线与五个等时姿态};
\end{scope}
\begin{scope}[xshift=2.1cm,yshift=-3.8cm]
\begin{scope}[scale=15]\KnifeOutline\end{scope}
\fill (0,0) circle(1.5pt);\node[below] at (0,-.25){$C$};\node[right] at (1.5,0){刀尖};
\draw[<->] (-2.25,-.65)--(1.5,-.65) node[midway,below]{$L=25\,\mathrm{cm}$};
\draw[<->] (0,.65)--(1.5,.65) node[midway,above]{$l=10\,\mathrm{cm}$};
\node at (0,-1.45){(b) 长度定义（局部放大）};
\end{scope}
\begin{scope}[xshift=8cm,yshift=-3.8cm]
\draw[kaxis] (-.3,0)--(2.5,0) node[right]{$x$};
\begin{scope}[rotate=60,scale=10]\KnifeOutline\end{scope}
\draw[vector] (0,0)--(30:2.1) node[above right]{$\vect v_C$};
\draw (.65,0) arc(0:30:.65);\node at (15:.95){$30^\circ$};
\draw (60:1.85) arc(60:360:1.85);\draw[->] (335:1.85) arc(335:360:1.85);
\node[left] at (-1.85,.45){$300^\circ$};
\node[right] at (.25,-.45){$C$};\draw[thin] (.25,-.45)--(0,0);
\node[above] at (60:2.05){$\theta=60^\circ$};
\node at (.3,-2.45){(c) 初姿态与最小正转角};
\end{scope}
\end{tikzpicture}\figcap{例2.96解图：尖端在刀轴正向；第5姿态刀尖恰达木靶，质心仍距靶 $l$，不是整段刀身穿靶。}\end{minipage}\end{center}}

\newcommand{\KnifeImpactDiagram}{%
\begin{center}\begin{minipage}{\linewidth}\centering\begin{tikzpicture}[>=Latex]
\begin{scope}[x=2cm,y=2cm]
\draw[kguide] (0,0)--(5,0);\draw[thick] (5,-.15)--(5,.5);
\begin{scope}[shift={(4.9,0)}]\KnifeOutline\end{scope}
\fill (0,0) circle(1.5pt);\node[above] at (0,0){$C_0$};\node[above] at (4.9,.15){$C_1$};
\draw[<->] (0,-.35)--(4.9,-.35) node[midway,below]{$d_1=s-l$};
\node at (2.5,.65){(a) 第一把刀击中靶点 $Q$};
\end{scope}
\begin{scope}[xshift=5cm,yshift=-3.9cm,x=15cm,y=15cm]
\begin{scope}[shift={(.10,0)}]\KnifeOutline\end{scope}
\begin{scope}[shift={(-.15,0)}]\KnifeOutline\end{scope}
\draw[thick] (.20,-.08)--(.20,.09);
\fill (-.05,0) circle(1.4pt) (-.15,0) circle(1.4pt) (.10,0) circle(1.4pt);
\node[above] at (-.05,.025){$P$};\node[below] at (-.15,-.015){$C_2$};\node[below] at (.10,-.015){$C_1$};\node[right] at (.20,0){$Q$};
\draw[<->] (-.05,.065)--(.20,.065) node[midway,above]{$L$};
\draw[<->] (-.15,-.085)--(-.05,-.085) node[midway,below]{$l$};
\node at (-.06,-.15){(b) 接触局部放大：$C_2P=l,\ PQ=L$};
\end{scope}
\end{tikzpicture}\figcap{例2.96碰撞几何：两刀仅在 $P$ 接触；第一把木柄远端就是第二把刀尖，故 $C_0C_2=s-L-l$。}\end{minipage}\end{center}}

\newcommand{\LightTimingDiagram}{%
\begin{center}\begin{minipage}{\linewidth}\centering\begin{tikzpicture}[>=Latex,scale=.95]
\coordinate (A) at (-3,0);\coordinate (B) at (3,0);\coordinate (O) at (-.7,2.5);
\draw[thick] (A)--(B);\draw[kguide] (A)--(O)--(B);
\foreach \p/\pos in {A/below left,B/below right,O/above}{\fill (\p) circle(1.5pt);\node[\pos] at (\p){$\p$};}
\node[above left] at (-1.95,1.2){$r_A$};\node[above right] at (1.4,1.2){$r_B$};
\draw[vector] (-.8,-.7)--(.8,-.7) node[midway,below]{运动员};
\node[above] at (0,0){$L=100\,\mathrm m$};
\end{tikzpicture}\figcap{例2.102解图：观察者 $O$ 接收起点 $A$、终点 $B$ 的光信号；传播路程一般不相等。}\end{minipage}\end{center}}

\newcommand{\ApparentTransverseDiagram}{%
\begin{center}\begin{minipage}{\linewidth}\centering\begin{tikzpicture}[>=Latex,scale=1]
\coordinate (E) at (0,0);\coordinate (P1) at (0,4);\coordinate (P2) at (2,2.6);\coordinate (F) at (0,2.6);
\draw[thick] (E)--(P1);\draw[thick] (E)--(P2);
\draw[vector] (P1)--(P2);\draw[kguide] (F)--(P2);
\fill (E) circle(1.6pt) (P1) circle(1.6pt) (P2) circle(1.6pt);
\node[below] at (E){地球};\node[above] at (P1){$P_1$};\node[right] at (P2){$P_2$};
\draw[<->] (-.9,0)--(-.9,4) node[midway,left]{$L$};
\node[above right] at (.85,3.55){$v\Delta t$};\node[below right] at (1.15,1.4){$D_2$};
\draw[kguide] (0,4.35)--(0,4.85) (P2)--(2,4.85);
\draw[<->] (0,4.85)--(2,4.85) node[midway,above]{$v\sin\theta\,\Delta t$};
\draw[<->] (-.4,2.6)--(-.4,4);\node[left,align=right] at (-1.3,3.3){接近量\\$v\cos\theta\,\Delta t$};\draw[thin] (-1.25,3.3)--(-.42,3.3);
\draw (0,3.35) arc(-90:-34.992:.65);\node at (.3,3.1){$\theta$};
\end{tikzpicture}\figcap{例2.103解图：横向与视线方向的位移分开标注；三角形给出 $D_2$，再对 $v\Delta t/L$ 作一阶展开。}\end{minipage}\end{center}}

\newcommand{\MarsRetrogradeDiagram}{%
\begin{center}\begin{minipage}{\linewidth}\centering\begin{tikzpicture}[>=Latex,scale=.9]
\foreach \t/\n/\dx/\dy in {-90/1/0/0,-60/2/3.5/0,-36.37/3/7/0,0/4/10.5/0,36.37/5/1.75/-4,60/6/5.25/-4,90/7/8.75/-4}{
\begin{scope}[shift={(\dx,\dy)}]
\draw[TextGray] (0,0) circle(.8);\draw[thick] (0,0) circle({.8*1.524});
\fill (0,0) circle(1.3pt);\node[above] at (0,0){$S$};
\pgfmathsetmacro{\ea}{.985609*\t-90}
\pgfmathsetmacro{\ma}{.985609*\t*pow(1.524,-1.5)-90}
\pgfmathsetmacro{\ex}{.8*cos(\ea)}\pgfmathsetmacro{\ey}{.8*sin(\ea)}
\pgfmathsetmacro{\mx}{.8*1.524*cos(\ma)}\pgfmathsetmacro{\my}{.8*1.524*sin(\ma)}
\draw[DeepBlue,dashed,thin,->] (\ex,\ey)--({\ex+1.7*(\mx-\ex)},{\ey+1.7*(\my-\ey)});
\fill (\ex,\ey) circle(1.5pt) (\mx,\my) circle(1.5pt);
\node[font=\scriptsize] at (\ea:.5){$E_{\n}$};
\ifdim\t pt<0pt \node[font=\scriptsize,left] at ({\mx-.13},{\my-.16}){$M_{\n}$};
\else \node[font=\scriptsize,right] at ({\mx+.13},{\my-.16}){$M_{\n}$};\fi
\node[font=\footnotesize] at (0,1.7){\n：$t=\t\,\mathrm d$};
\end{scope}}
\node at (5.25,-6){虚线由地球经过火星，指向背景星空；各小图比例相同};
\begin{scope}[xshift=5.25cm,yshift=-8.6cm]
\draw[kaxis] (-3.7,0)--(3.8,0) node[right]{$t/\mathrm d$};
\draw[kaxis] (0,-1.9)--(0,1.45) node[above]{$\lambda$};
\MarsCoplanarPath
\foreach \t/\n/\pos in {-90/1/above left,-60/2/above left,-36.37/3/above,0/4/below left,36.37/5/below,60/6/below right,90/7/below right}{
\pgfmathsetmacro{\ll}{atan2(1.524*sin(.985609*\t*pow(1.524,-1.5))-sin(.985609*\t),1.524*cos(.985609*\t*pow(1.524,-1.5))-cos(.985609*\t))/12}
\fill (\t/27,\ll) circle(1.5pt);\node[\pos] at (\t/27,\ll){\n};}
\draw[<->] (-1.347,-1.75)--(1.347,-1.75) node[midway,below]{逆行};
\node at (0,2.15){同平面模型：七个时刻的视经度};
\node at (0,-2.65){第3、5点：$\dot\lambda=0$};
\end{scope}
\end{tikzpicture}\figcap{例2.105解图：七组位置与七个视方向一一对应；冲日取 $t=0$，第3至第5点之间视经度随时间减小。}\end{minipage}\end{center}}

\newcommand{\MarsInclinedDiagram}{%
\begin{center}\begin{minipage}{\linewidth}\centering\begin{tikzpicture}[>=Latex]
\draw[kaxis] (-4.4,-2.1)--(4.6,-2.1) node[right]{视经度增加};
\draw[kaxis] (-4.4,-2.1)--(-4.4,3.05) node[above]{视纬度增加};
\MarsInclinedPath
\node[above] at (0,3.25){三维圆轨道示例：轨道倾角 $1.85^\circ$};
\node at (0,-2.8){箭头表示时间先后；蓝色区间才是逆行};
\end{tikzpicture}\figcap{回环与计算模型的区别：加入非零倾角后视纬度也会变化；蓝段按视经度减小的条件确定，不把整个回环都称为逆行。}\end{minipage}\end{center}}

\newcommand{\RomerDiagram}{%
\begin{center}\begin{minipage}{\linewidth}\centering\begin{tikzpicture}[>=Latex,scale=1]
\coordinate (S) at (0,0);\coordinate (J) at (5.2,0);
\pgfmathsetmacro{\ang}{acos(1/5.2)}
\draw[thick] (S) circle(1);\draw[TextGray] ({5.2*cos(-10)},{5.2*sin(-10)}) arc(-10:10:5.2);
\foreach \a/\lab/\pos in {\ang/E_+/above left,-\ang/E_-/below left}{
\coordinate (E) at (\a:1);\draw[kguide] (S)--(E);\draw[thick] (E)--(J);\fill (E) circle(1.6pt);\node[\pos] at (E){$\lab$};
\draw[vector] (E)--++({-.65*sin(\a)},{.65*cos(\a)});}
\fill (S) circle(1.6pt) (J) circle(1.6pt);\node[left] at (S){$S$};\node[right] at (J){$J$};
\draw[kguide] (S)--(J);
\draw[<->] (0,-2.25)--(5.2,-2.25) node[midway,below]{$R_J=5.2R_E$};
\node at (2,1.65){$E_+$：距离增加，周期最长};\node at (2,-1.65){$E_-$：距离减小，周期最短};
\begin{scope}[xshift=9cm]
\draw (0,0) circle(.65);\fill (0,0) circle(2pt);\node[below] at (0,0){$J$};
\fill (45:.65) circle(1.6pt);\node[above right] at (45:.65){木卫I};
\draw[->] (50:.9) arc(50:170:.9);\node at (0,-1.2){卫星轨道放大};
\end{scope}
\end{tikzpicture}\figcap{例2.106解图：$SE_\pm\perp E_\pm J$，两切点为周期极值位置；木卫I单独放大，不与日心轨道混用比例。}\end{minipage}\end{center}}

\newcommand{\AimLocusDiagram}{%
\begin{center}\begin{minipage}{\linewidth}\centering\begin{tikzpicture}[>=Latex,scale=.9]
\begin{scope}
\draw[thick] (-.4,0)--(4.1,0);\draw[thick] (3.6,0)--(3.6,4.4);
\coordinate (O) at (0,.6);\coordinate (Q) at (3.6,3.7);\coordinate (P) at (3.6,1.6);
\draw[kguide] (O)--(Q);\draw[kcurve,domain=0:1,samples=70,variable=\t] plot({3.6*\t},{.6+3.1*\t-2.1*\t^2});
\foreach \p in {O,Q,P}{\fill (\p) circle(1.5pt);}
\node[above left=3pt] at (O){$O$};\node[right] at (Q){$Q$};\node[below left] at (3.5,1.45){$P\ (z=h)$};
\draw[<->] (0,-.5)--(3.6,-.5) node[midway,below]{$r$};
\draw[<->] (4.65,.6)--(4.65,3.7) node[midway,right]{$Z$};
\draw[kguide] (0,.6)--(4.65,.6);
\node[left] at (0,.25){$z_0$};
\node at (1.8,5){(a) 包含 $OQ$ 的竖直截面};
\end{scope}
\begin{scope}[xshift=9cm]
\draw[kaxis] (-2.4,0)--(2.6,0) node[right]{$x$};\draw[kaxis] (0,0)--(0,4.6) node[above]{$z$};
\draw[DeepBlue,thick] (0,2.7) circle(1.5);\draw[thick] (-2.4,.7)--(2.4,.7);
\coordinate (Q) at (1.2,3.6);\coordinate (P) at (1.2,.7);
\draw[kguide] (Q)--(P);\fill (Q) circle(1.5pt) (P) circle(1.5pt) (0,2.7) circle(1.5pt);
\node[above right] at (Q){$Q$};\node[below right] at (P){$P$};
\node[right] at (2.45,.7){$h$};
\node[left] at (-.1,2.7){$K$};
\node at (0,5.25){(b) 正视墙面};
\end{scope}
\end{tikzpicture}\figcap{例2.109解图：瞄准点 $Q$ 位于以 $K$ 为圆心的墙面圆上，实际命中点 $P$ 位于水平线；两图分别说明投掷平面与墙面轨迹。}\end{minipage}\end{center}}
